% =====================================================================
%  exam-frequency-report / assets/report-preamble.tex
%  Reusable preamble for an exam-frequency report.
%
%  Compile with XeLaTeX (twice).  Load-bearing choices, each of which
%  fixes a real failure mode -- see the skill's Pitfalls table:
%    * ctexart + xelatex          -> CJK actually renders
%    * amsmath only, before newtxmath -> avoids the \Bbbk clash with amssymb
%    * xcolor[table]              -> \rowcolor is defined
%    * \noindent before \rule/\tabularx -> no 17pt \parindent overfull
%  Adjust \setReportMeta and the geometry block per task.
% =====================================================================
\documentclass[11pt,fontset=fandol]{ctexart}

% --- page geometry ---------------------------------------------------
% Keep the long edge at a size you have already validated (A4 height
% preserves the vertical layout).  To match a reading device, set
%   paperwidth = paperheight * Wpx / Hpx
% e.g. 13-inch iPad Air M2 portrait (2048x2732):
%   paperheight=297mm, paperwidth=297*512/683=222.6413mm
\usepackage[paperwidth=222.6413mm,paperheight=297mm,
            left=1.9cm,right=1.9cm,top=2.1cm,bottom=1.7cm,headsep=0.5cm]{geometry}

% --- math: order matters ---------------------------------------------
% DO NOT load newtxtext in a CJK (ctex) document.
% On macOS ctex falls back to the macold fontset, whose \ttfamily CJK slot
% points at STFangsong -- absent from current macOS.  Any Chinese inside
% \texttt (e.g. a path written as \code{试卷/教材}) then fails with
%   ! Font TU/STFangsong ... not loadable
% and those glyphs come out as mojibake.  Verified: newtxtext triggers it,
% newtxmath is harmless, and \setCJKmonofont does NOT repair it.
% Pure-ASCII (non-ctex) documents may load newtxtext normally.
\usepackage{amsmath}
\usepackage{newtxmath}

% --- tables, colour, boxes -------------------------------------------
\usepackage[table]{xcolor}
\usepackage{booktabs}
\usepackage{longtable}
\usepackage{tabularx}
\usepackage{array}
\usepackage{multirow}
\usepackage{makecell}
\usepackage{ragged2e}
\usepackage{enumitem}
\usepackage{tcolorbox}
\tcbuselibrary{skins,breakable}
\usepackage{pgfplots}
\pgfplotsset{compat=1.18}
\usetikzlibrary{patterns}
\usepackage{titlesec}
\usepackage{fancyhdr}
\usepackage{caption}
\usepackage[hidelinks,bookmarksnumbered]{hyperref}

% --- palette ---------------------------------------------------------
\definecolor{ink}{HTML}{1F2933}
\definecolor{deepblue}{HTML}{1B3A5C}
\definecolor{accentA}{HTML}{B23A48}   % S tier
\definecolor{accentB}{HTML}{D97706}   % A tier
\definecolor{accentC}{HTML}{2F6F4E}   % B tier
\definecolor{accentD}{HTML}{6B7280}   % C tier
\definecolor{tintA}{HTML}{FBE9EB}
\definecolor{tintB}{HTML}{FDF2E3}
\definecolor{tintC}{HTML}{E8F2EC}
\definecolor{tintD}{HTML}{F0F1F3}
\definecolor{rulegray}{HTML}{C9CED6}
\definecolor{tabhead}{HTML}{1B3A5C}

% --- headings --------------------------------------------------------
\titleformat{\section}{\normalfont\Large\bfseries\color{deepblue}}{\thesection}{0.7em}{}
\titleformat{\subsection}{\normalfont\large\bfseries\color{deepblue!85!black}}{\thesubsection}{0.6em}{}
\titlespacing*{\section}{0pt}{1.6em}{0.7em}
\titlespacing*{\subsection}{0pt}{1.2em}{0.5em}

% --- running head ----------------------------------------------------
\pagestyle{fancy}
\fancyhf{}
\renewcommand{\headrulewidth}{0.4pt}
\fancyhead[L]{\footnotesize\color{deepblue}\ReportShortTitle}
\fancyhead[R]{\footnotesize\color{deepblue}\thepage}
\fancyfoot[C]{\scriptsize\color{accentD}\ReportSourceLine}
\newcommand{\ReportShortTitle}{Exam-frequency report}
\newcommand{\ReportSourceLine}{}

% --- helpers ---------------------------------------------------------
% Section-code tag.  Use the textbook's OWN numbering.
\newcommand{\secode}[1]{\texttt{\color{deepblue}#1}}
\newcommand{\tierS}{\textbf{\color{accentA}S}}
\newcommand{\tierA}{\textbf{\color{accentB}A}}
\newcommand{\tierB}{\textbf{\color{accentC}B}}
\newcommand{\tierC}{\textbf{\color{accentD}C}}
\newcommand{\hot}[1]{\textcolor{accentA}{\bfseries #1}}
\newcommand{\code}[1]{\texttt{\small #1}}
\newcommand{\thd}[1]{\textcolor{white}{\bfseries #1}}
\newcolumntype{L}[1]{>{\RaggedRight\arraybackslash}p{#1}}
\newcolumntype{C}[1]{>{\Centering\arraybackslash}p{#1}}
% NOTE: \newlist needs an explicit label= or enumitem raises "Undefined label".
\newlist{qparts}{enumerate}{1}
\setlist[qparts]{label=(\alph*),leftmargin=1.5em,itemsep=2.5pt,topsep=3pt,parsep=0pt}
\newlist{qnum}{enumerate}{1}
\setlist[qnum]{label=(\arabic*),leftmargin=1.5em,itemsep=2.5pt,topsep=3pt,parsep=0pt}
\newlist{bul}{itemize}{1}
\setlist[bul]{label=\textbullet,leftmargin=1.6em,itemsep=1.5pt,topsep=2pt,parsep=0pt}
\captionsetup{font={small,color=deepblue},labelfont=bf,skip=5pt}

% --- callout boxes ---------------------------------------------------
\newtcolorbox{keybox}[1][]{colback=tintC,colframe=accentC,boxrule=0.7pt,arc=2pt,
  left=8pt,right=8pt,top=6pt,bottom=6pt,fonttitle=\bfseries\color{white},
  coltitle=white,colbacktitle=accentC,breakable,#1}
\newtcolorbox{warnbox}[1][]{colback=tintA,colframe=accentA,boxrule=0.7pt,arc=2pt,
  left=8pt,right=8pt,top=6pt,bottom=6pt,fonttitle=\bfseries\color{white},
  coltitle=white,colbacktitle=accentA,breakable,#1}
\newtcolorbox{infobox}[1][]{colback=tintB,colframe=accentB,boxrule=0.7pt,arc=2pt,
  left=8pt,right=8pt,top=6pt,bottom=6pt,fonttitle=\bfseries\color{white},
  coltitle=white,colbacktitle=accentB,breakable,#1}

% --- one-item-per-page layout (review sets) --------------------------
% tcolorbox `height fill` absorbs the remaining page height.
% Do NOT combine with `breakable`.
\newtcolorbox{workspace}{height fill,enhanced,colback=white,colframe=gray!45,
  boxrule=0.5pt,arc=1.5pt,left=9pt,right=9pt,top=6pt,bottom=6pt}

% --- pgfplots reference lines on a symbolic axis ---------------------
% A numeric line on `symbolic x coords` fails with
%   "the input coordinate '...' has not been defined".
% Draw it in axis-description coordinates instead, e.g. y=33.75 of ymax=60:
%   \draw[accentA,dashed,thick]
%     (axis description cs:0,0.5625) -- (axis description cs:1,0.5625);
%   \node[anchor=south west,font=\scriptsize,color=accentA]
%     at (axis description cs:0.02,0.578) {four-year mean};


\setlength{\parindent}{0pt}\setlength{\parskip}{5pt}
\setlength{\headheight}{14pt}
\newcommand{\E}{\mathbb E}\newcommand{\Var}{\operatorname{Var}}\newcommand{\Cov}{\operatorname{Cov}}
\newcommand{\key}[1]{\par\textbf{\color{deepblue}#1}\par}
\newcommand{\pagehead}[1]{\clearpage\section*{#1}}
\renewcommand{\ReportShortTitle}{随机信号分析 · 第二章考频分析}
\renewcommand{\ReportSourceLine}{四套真题 · 2022--23 至 2025--26 · 2026-10-13}
\begin{document}

\thispagestyle{empty}\vspace*{15mm}
{\Huge\bfseries\color{deepblue}第二章：随机信号\par 历年考频分析报告\par}
\vspace{8mm}{\Large 真题审计 · 作业覆盖 · 七天复习路线\par}
\vspace{12mm}
\begin{keybox}[title=先看结论]
沿用第一章的唯一主考点归属后，第二章四套试卷的核心分值为
\[\boxed{29\;/\;0\;/\;7\;/\;0},\qquad\text{均值 }9.00\text{ 分/套}.\]
核心共 36 分，占四套总分的 9\%；2/4 套有核心小问，0/4 套有完整 25 分大题全部归本章。
但高斯信号在四套试卷中均被用到。\textbf{核心分少不代表可以跳过本章：很多分数已经归入第一章，或以平稳性、滤波、带通题形式出现。}
\end{keybox}
\vspace{5mm}
\textbf{复习顺序：}先练高斯有限维分布与四阶矩，再练互相关、随机相位与独立性，最后查漏样本函数和离散分布。
\vspace{6mm}
\begin{infobox}[title=本次配套]
英文复习题集 21 题，7 天、每天 3 题；中文详解逐题对应。第二次作业 10 题全部覆盖，另加 11 题真题改编或补缺训练。
作业必做：2.3、2.13、2.16、2.19；选做：2.1、2.4、2.7、2.8、2.11、\textbf{2.18}。
\end{infobox}
\vfill
教材：李晓峰等《随机信号分析（第5版）》第二章，书 pp.37--62。\\
统计样本：工作区四套真题；不是考试预测，也不是官方评分细则。
\pagehead{1\quad 统计口径与跨章边界}
\key{核心与间接}
每个小问只保留一个主考点，全部分数计给它。间接标签只统计小问数和命中卷数，绝不累加分值。
本报告保留第一章已登记的全部归属；第一章和第二章不能重复拿同一道小问的分数。
\begin{center}\begin{tabular}{lrrrr}\toprule
学年 & 2022--23 & 2023--24 & 2024--25 & 2025--26\\\midrule
试卷总分 &100&100&100&100\\
第一章（已锁定）&29&33&24&49\\
第二章（本报告）&29&0&7&0\\
留给其余章节&42&67&69&51\\\bottomrule\end{tabular}\end{center}
\key{三个需要解释的归属}
\begin{bul}
\item 2022--23 Q3(b) 的 15 分归 §2.4.3：由(a)的功率谱得到相关函数，平方器要求计算高斯信号四阶矩，对应教材习题2.17。滤波环节(a)仍归第五章；本小问只计一次。
\item 2023--24 Q1(c) 归 §3.4.2：题目只给功率谱，先做傅里叶反变换才能开始数值计算；§2.3.2 记为间接涉及。
\item 2024--25 Q3(c) 归 §5.1.2：必须先由RC两支路频响确定互相关；第二章的相关展开是工具，不重复计分。
\end{bul}
\key{可复算的检验}
沿用清单中已核验的100分结构；逐小问登记后，每卷总分仍为100。
\[6+8+15+7=36=29+0+7+0.\]
各小节核心分之和等于各试卷本章分之和。数据表另存为同目录 JSON，便于后续章节延用。
\begin{warnbox}[title=解释范围]
如果把“出现在随机过程题里”的所有概率运算都算第二章，数字会明显变大，却与第一章重复。本报告的9分是\textbf{固定跨章归属后的独占均值}，不是第二章知识实际支持的总分。
\end{warnbox}
\pagehead{2\quad 教材知识地图}
\small\begin{tabularx}{\textwidth}{@{}l L{5cm} X@{}}\toprule 小节 & 教材标题 & 本次训练重点\\\midrule
2.1.1 & 概念与定义 & 样本函数与固定时刻的随机变量\\
2.1.2 & 基本概率特性 & 离散原子、边缘分布和归一化\\
2.1.3 & 基本数字特征 & 均值、相关、协方差的区分\\
2.2.1 & 随机正弦信号 & 幅度与相位随机性\\
2.2.2 & 伯努利随机序列 & 独立重采样与共享同一次试验\\
2.2.3 & 半随机二进制传输信号 & 本样本未单独考查\\
2.2.4 & 泊松过程 & 本样本未单独考查\\
2.3.1 & n阶概率特性 & 所有有限维联合分布\\
2.3.2 & 联合特性 & 互相关及正交、无关\\
2.3.3 & 相关函数与协方差函数的性质 & 半正定、对称性与线性展开\\
2.3.4 & 微分和积分（选学） & 均方微积分；积分交换条件\\
2.4.1 & 多维高斯分布 & 先检验协方差是否奇异\\
2.4.2 & 高斯随机变量的性质 & 线性变换、联合高斯与独立性\\
2.4.3 & 高斯随机信号 & 有限维分布、高阶矩、线性系统\\
2.5 & 独立信号 & 所有有限采样向量之间的独立性\\
\bottomrule\end{tabularx}\normalsize
\vspace{5mm}\begin{infobox}[title=章节界线]
平稳性和功率谱属于第三章，各态历经性属于第四章，LTI滤波属于第五章，带通正交分量属于第六章。第二章提供分布与矩的基础。
本书第二章到 §2.5 为止，\textbf{没有 §2.6}。§2.3.4 为选学，本次仍列入并注明零核心。
\end{infobox}
\pagehead{3\quad 2022--23：逐小问审计}
\small\renewcommand{\arraystretch}{1.5}\begin{tabularx}{\textwidth}{@{}l r l L{3.5cm} X@{}}\toprule 小问 & 分 & 唯一主考点 & 第二章间接标签 & 判定依据\\\midrule
Q1a & 6 & \textbf{2.1.1} & — & 固定随机参数画样本函数\\
Q1b & 9 & 1.5.1 & 2.1.2 & 平移均匀变量的特征函数\\
Q1c & 10 & 1.4.2 & 2.1.3 & 随机直线的均值与协方差\\
Q2a & 8 & \textbf{2.2.1} & 2.5 & 均匀随机相位平均为零\\
Q2b & 8 & 3.1.1 & 2.2.1 / 2.5 & 计算相关后判断广义平稳\\
Q2c & 9 & 3.4.2 & — & 由相关函数变换得到调制谱\\
Q3a & 10 & 5.2.1 & — & 白噪声通过RC低通\\
Q3b & 15 & \textbf{2.4.3} & 2.4.2 & 承接(a)，平方器输出需高斯四阶矩\\
Q4a & 15 & 5.1.2 & — & LTI输出均值、谱和相关\\
Q4b & 10 & 1.3.1 & 2.4.3 / 2.4.1 & 高斯输入经线性系统保持高斯\\
\bottomrule\end{tabularx}\normalsize\vspace{5mm}
本章核心：29 分。粗体主考点属于第二章；第一章归属与既有报告一致。
\begin{infobox}[title=读题核对记录]Q3(b)：$N(t)$ 明确为高斯白噪声，$Y(t)=aX^2(t)$；不能把 $Y$ 也当成高斯。$R_Y$ 中必须保留均值平方项。\end{infobox}\pagehead{4\quad 2023--24：逐小问审计}
\small\renewcommand{\arraystretch}{1.5}\begin{tabularx}{\textwidth}{@{}l r l L{3.5cm} X@{}}\toprule 小问 & 分 & 唯一主考点 & 第二章间接标签 & 判定依据\\\midrule
Q1a & 12 & 3.1.2 & 2.3.2 / 2.5 & 和差信号的联合平稳性\\
Q1b & 6 & 3.4.2 & — & 和信号的功率谱\\
Q1c & 7 & 3.4.2 & 2.3.2 & 由已知谱反变换并组合互相关\\
Q2a & 10 & 5.2.1 & — & 白噪声滤波后的相关函数\\
Q2b & 7 & 5.1.2 & 2.3.3 & 延时差分的输出谱\\
Q2c & 8 & 1.6 & 2.4.3 & 白噪声未明说高斯，结论需条件\\
Q3a & 4 & 1.2.1 & — & 联合密度归一化\\
Q3b & 8 & 1.4.1 & — & 由密度求均值\\
Q3c & 6 & 1.4.2 & — & 由密度求方差\\
Q3d & 7 & 1.4.2 & — & 联合矩与相关系数\\
Q4a & 4 & 6.3.1 & — & 希尔伯特变换与正交分量\\
Q4b & 15 & 6.3.2 & — & 正交分量功率谱与互谱\\
Q4c & 6 & 6.3.2 & — & 载频位置改变互谱\\
\bottomrule\end{tabularx}\normalsize\vspace{5mm}
本章核心：0 分。粗体主考点属于第二章；第一章归属与既有报告一致。
\begin{infobox}[title=读题核对记录]Q2 只写 white noise，没有给 Gaussian。Q2(c) 的“是否高斯”不能仅由白噪声性质肯定；详解第21题给出条件式答案。\end{infobox}\pagehead{5\quad 2024--25：逐小问审计}
\small\renewcommand{\arraystretch}{1.5}\begin{tabularx}{\textwidth}{@{}l r l L{3.5cm} X@{}}\toprule 小问 & 分 & 唯一主考点 & 第二章间接标签 & 判定依据\\\midrule
Q1a & 7 & \textbf{2.3.2} & 2.2.1 & 离散正弦和余弦的互相关\\
Q1b & 9 & 1.4.3 & 2.3.2 / 2.5 & 正交、无关与独立性区别\\
Q1c & 9 & 3.1.2 & 2.3.2 & 联合平稳性判断\\
Q2a & 12 & 3.1.1 & 2.3.2 / 2.2.1 & 正交调制和信号的平稳条件\\
Q2b & 8 & 3.4.2 & — & 由平稳条件求功率谱\\
Q2c & 5 & 1.4.3 & 2.3.2 & 无关条件下交叉项简化\\
Q3a & 9 & 5.2.1 & — & RC电阻、电容支路频响\\
Q3b & 6 & 5.2.1 & — & 由支路谱求相关\\
Q3c & 10 & 5.1.2 & 2.3.2 & 先由支路频响求互相关\\
Q4a & 9 & 6.3.2 & — & 理想带通噪声及正交分量\\
Q4b & 6 & 6.3.2 & — & 同相与正交分量互谱\\
Q4c & 10 & 1.6 & 2.4.1 / 2.4.2 / 2.4.3 & 联合高斯密度及无关推独立\\
\bottomrule\end{tabularx}\normalsize\vspace{5mm}
本章核心：7 分。粗体主考点属于第二章；第一章归属与既有报告一致。
\begin{infobox}[title=读题核对记录]Q1 的 $n$ 为整数，故 $\sin(2\pi n+\phi)=\sin\phi$、$\cos(2\pi n+\phi)=\cos\phi$。卷面日期空白，学年标签沿用文件清单。\end{infobox}\pagehead{6\quad 2025--26：逐小问审计}
\small\renewcommand{\arraystretch}{1.5}\begin{tabularx}{\textwidth}{@{}l r l L{3.5cm} X@{}}\toprule 小问 & 分 & 唯一主考点 & 第二章间接标签 & 判定依据\\\midrule
Q1a & 6 & 1.1.2 & 2.1.2 & 同一次投币定义的单时刻密度\\
Q1b & 6 & 1.2.1 & 2.3.1 & 同一次投币定义的二时刻密度\\
Q1c & 6 & 1.4.2 & 2.1.3 & 两个样本函数的协方差\\
Q1d & 7 & 1.2.1 & 2.2.2 / 2.5 & 每个时刻重新独立投币\\
Q2a & 9 & 1.5.1 & 2.2.1 / 2.5 & 相位特征函数控制平稳条件\\
Q2b & 7 & 3.4.2 & — & 随机幅度调制的功率谱\\
Q2c & 9 & 4.1.1 & 2.1.3 & 利用时间平均方差判断均值各态历经\\
Q3a & 10 & 5.1.2 & 2.3.3 & 两径LTI系统输出相关与谱\\
Q3b & 10 & 5.4.3 & — & 匹配滤波后的相关与谱\\
Q3c & 5 & 5.3.2 & — & 等效矩形带宽与时延\\
Q4a & 10 & 6.3.2 & — & 非中心载频的正交分量谱\\
Q4b & 7 & 1.6 & 2.4.1 / 2.4.2 / 2.4.3 & 同一时刻二维高斯分布\\
Q4c & 8 & 1.4.2 & 2.4.1 / 2.4.3 & 异时刻协方差决定联合密度\\
\bottomrule\end{tabularx}\normalsize\vspace{5mm}
本章核心：0 分。粗体主考点属于第二章；第一章归属与既有报告一致。
\begin{infobox}[title=读题核对记录]Q1 中必须保留 $t^2,n^2$。Q2(c) 原文为 $R_A(\tau)=e^{-|\tau|}$，不是 $R_X$；它与前述 $m_A\ne0$ 在通常有限二阶矩WSS模型下不相容，因此本题未纳入第二章训练。Q4 图中 $144\pi\le|\omega|\le156\pi$，高度1，载频 $148\pi$ 不在频带中心。\end{infobox}\pagehead{7\quad 小节考频：分值与涉及率}\small\renewcommand{\arraystretch}{1.3}\begin{tabularx}{\textwidth}{@{}l X r r r r r@{}}\toprule 小节 & 标题 & 核心分 & 核心卷 & 间接问 & 间接卷 & 并集卷\\\midrule2.1.1 & 概念与定义 & 6 & 1/4 & 0 & 0/4 & 1/4\\
2.1.2 & 基本概率特性 & 0 & 0/4 & 2 & 2/4 & 2/4\\
2.1.3 & 基本数字特征 & 0 & 0/4 & 3 & 2/4 & 2/4\\
2.2.1 & 随机正弦信号 & 8 & 1/4 & 4 & 3/4 & 3/4\\
2.2.2 & 伯努利随机序列 & 0 & 0/4 & 1 & 1/4 & 1/4\\
2.2.3 & 半随机二进制传输信号 & 0 & 0/4 & 0 & 0/4 & 0/4\\
2.2.4 & 泊松过程 & 0 & 0/4 & 0 & 0/4 & 0/4\\
2.3.1 & n阶概率特性 & 0 & 0/4 & 1 & 1/4 & 1/4\\
2.3.2 & 联合特性 & 7 & 1/4 & 7 & 2/4 & 2/4\\
2.3.3 & 相关函数与协方差函数的性质 & 0 & 0/4 & 2 & 2/4 & 2/4\\
2.3.4 & 微分和积分（选学） & 0 & 0/4 & 0 & 0/4 & 0/4\\
2.4.1 & 多维高斯分布 & 0 & 0/4 & 4 & 3/4 & 3/4\\
2.4.2 & 高斯随机变量的性质 & 0 & 0/4 & 3 & 3/4 & 3/4\\
2.4.3 & 高斯随机信号 & 15 & 1/4 & 5 & 4/4 & 4/4\\
2.5 & 独立信号 & 0 & 0/4 & 6 & 4/4 & 4/4\\
\bottomrule\end{tabularx}\normalsize
\vspace{4mm}
“间接问”是一小问对一个小节至多计一次；同一小问可以涉及多个小节，故该列不能相加当作题数或分值。
“并集卷”是核心或间接涉及的不同试卷数；0/4、1/4、2/4、3/4、4/4分别对应0\%、25\%、50\%、75\%、100\%。
\begin{keybox}[title=读表结论]
§2.4.3 四套均涉及，是最值得先复习的知识连接点。§2.5 虽然核心为0，也在四套中作为独立性前提出现。
§2.4.1 的多维高斯分布没有本章独占分数，却支撑后续联合密度题，不能据“0分”略过。
\end{keybox}
\pagehead{8\quad 趋势与核心分排名}
\begin{center}\begin{tikzpicture}\begin{axis}[width=.86\textwidth,height=6cm,ybar,bar width=24pt,ymin=0,ymax=35,ytick={0,7,14,21,29,35},symbolic x coords={22--23,23--24,24--25,25--26},xtick=data,ylabel={核心分},nodes near coords,axis lines*=left]
\addplot[fill=deepblue!65] coordinates {(22--23,29)(23--24,0)(24--25,7)(25--26,0)};
\draw[accentA,dashed] (axis description cs:0,0.25714)--(axis description cs:1,0.25714);
\node[anchor=west,font=\small,color=accentA] at(axis description cs:.65,.31){均值9分};
\end{axis}\end{tikzpicture}\end{center}
\begin{center}\begin{tikzpicture}\begin{axis}[width=.86\textwidth,height=6cm,xbar,bar width=18pt,xmin=0,xmax=18,symbolic y coords={2.1.1,2.3.2,2.2.1,2.4.3},ytick=data,xlabel={四套核心分合计},nodes near coords,axis lines*=left]
\addplot[fill=accentC!70] coordinates {(6,2.1.1)(7,2.3.2)(8,2.2.1)(15,2.4.3)};
\end{axis}\end{tikzpicture}\end{center}
最近一套核心为0，是归属边界造成的读数：其Q1仍完整训练过程的有限维分布，Q4仍用高斯信号工具。四个样本不支持“第二章越来越不考”的趋势判断。
\pagehead{9\quad 复习等级与补缺优先级}
沿用第一章按核心分、命中卷数判断 S/A/B/C，并将前提型单列。既有资料只明确 S 为核心累计至少25分，未给其余档的统一数值切线；本章沿用已完成报告的实例尺度，\textbf{不补造精确阈值}。
\begin{center}\begin{tabularx}{\textwidth}{@{}l L{4.6cm} X@{}}\toprule 等级 & 考点 & 解释与行动\\\midrule
S & 本章无 & 最高15分，未达既有25分标准。\\
A & §2.4.3（15分，1/4核心） & 四套均涉及；先掌握有限维高斯和四阶矩。\\
B & §2.2.1（8）、§2.3.2（7）、§2.1.1（6） & 均为1/4核心；分别练相位平均、互相关展开、样本函数。\\
前提型 & §2.5、§2.4.1、§2.4.2 & 无独占分但反复涉及；优先级不按零分处理。\\
C & 其余零核心小节 & 结合作业保底；零命中不是不考承诺。\\\bottomrule\end{tabularx}\end{center}
\key{三个作业之外的高价值补缺}
\begin{qnum}
\item \textbf{高斯四阶矩：}作业未列2.17，2022--23 Q3(b)却要求它。配套第3、4题给证明与平方器应用。
\item \textbf{高斯前提检查：}白噪声不等于高斯白噪声；边缘高斯不等于联合高斯。配套第1、21题专练。
\item \textbf{相关性质和独立信号：}会写矩的定义不足以判断过程独立；必须考虑任意有限采样向量。配套第10、12题补齐。
\end{qnum}
\key{两道不能套公式的已布置题}
2.11只给“两两独立”，不足以分解三因子的期望；2.18的三维协方差矩阵奇异，普通三维PDF不存在。题面照录，答案中明确解释，不悄悄修改条件。
\pagehead{10\quad 第二次作业覆盖审计}
\small\begin{tabularx}{\textwidth}{@{}l l L{3cm} X l@{}}\toprule 题号 & 性质 & 小节 & 对应能力／缺口 & 复习题\\\midrule
2.3 & 必做 & 2.1.2 / 2.3.1 & 同次投币的边缘与联合CDF；接最新Q1 & 13\\
2.13 & 必做 & 2.2.1 / 2.3.4 & 随机相位、微分与电流均方值 & 19\\
2.16 & 必做 & 2.4.3 / 2.4.1 & 有限维高斯、退化采样时刻 & 5\\
2.19 & 必做 & 2.4.1 & 多维PDF、逆矩阵及联合特征函数 & 2\\
2.1 & 选做 & 2.1.1 / 2.2.1 & 样本函数、缩放密度、零振幅原子 & 15\\
2.4 & 选做 & 2.2.2 & 独立比特串与平均向量 & 16\\
2.7 & 选做 & 2.4.2 / 2.1.3 & 高斯线性组合与时变均值方差 & 18\\
2.8 & 选做 & 2.2.2 / 2.1.3 & 离散密度、对角协方差 & 17\\
2.11 & 选做 & 2.3.2 / 2.5 & 互相关与两两独立不足的反例 & 9\\
2.18 & 选做 & 2.4.1 / 2.4.3 & 奇异协方差，密度不存在的判断 & 6\\\bottomrule\end{tabularx}\normalsize
\vspace{6mm}
\begin{keybox}[title=覆盖结果]
必做4/4、选做6/6、全部10/10均在复习题集中逐字沿用既有英文题面。高斯相关题放在前两天，随后衔接随机相位与互相关。
作业训练覆盖强，但没有直接训练2.17四阶矩证明与白噪声高斯性判断，所以增加对应真题改编。
\end{keybox}
\key{暂未展开}
§2.2.3半随机二进制传输信号、§2.2.4泊松过程在四套样本中无本章核心或必要间接标签，且不在本次作业清单，21题未另设专项。
选学§2.3.4保留作业2.13，并增补2.14积分，避免仅会微分不会处理相关函数的双重积分。
\pagehead{11\quad 七天路线与使用方法}
\begin{center}\begin{tabularx}{\textwidth}{@{}c l X@{}}\toprule 天 & 题号 & 主题与完成标准\\\midrule
1&1--3&高斯定义、2.19、2.17：会判断PDF存在，会证明四阶矩。\\
2&4--6&平方器真题、2.16、2.18：检查非线性与奇异情形。\\
3&7--9&随机相位真题、离散正余弦真题、2.11：先查独立性条件。\\
4&10--12&相关性质、随机直线真题、独立信号：补概念与展开步骤。\\
5&13--15&2.3、最新投币真题、2.1：共享一次随机试验与重新采样。\\
6&16--18&2.4、2.8、2.7：离散序列与高斯线性组合。\\
7&19--21&2.13、2.14、白噪声真题辨析：微积分与条件检查。\\\bottomrule\end{tabularx}\end{center}
每天先独立完成3题，再核对中文详解。错题标记为：条件未读清、模型建立错误、积分/代数错误、退化情形遗漏；次日先重做错题关键步骤。
\begin{infobox}[title=与提交日期的关系]
作业布置日期为2026-10-13，截止2026-10-27。七天计划是复习节奏，不新增截止日期；可在余下时间复做错题与完成提交稿。
\end{infobox}
\key{复习题来源标记}
“Textbook”题目来自教材重排；本次作业题沿用既有英译，不称官方英译。
“Adapted from”表示改编或提取考点，新增条件明确写出，不冒充原卷全文。
“Recall / extension”是自编的概念核验，无真题分值。
\pagehead{12\quad 公式速查：先条件，后代入}
\key{均值、相关与协方差（实值过程）}
\[m_X(t)=\E X(t),\quad R_X(s,t)=\E[X(s)X(t)],\quad C_X(s,t)=R_X(s,t)-m_X(s)m_X(t).\]
\[R_{XY}(s,t)=\E[X(s)Y(t)],\quad C_{XY}=R_{XY}-m_X(s)m_Y(t).\]
\[C_X(t,t)=\Var X(t),\quad R_X(t,t)=\E X^2(t).\]
\key{高斯向量：$C$须半正定，普通PDF须正定}
\[\phi_{\boldsymbol X}(\boldsymbol u)=\exp\{j\boldsymbol u^T\boldsymbol\mu-\tfrac12\boldsymbol u^TC\boldsymbol u\}.\]
\[f(\boldsymbol x)=\frac{\exp[-\tfrac12(\boldsymbol x-\boldsymbol\mu)^TC^{-1}(\boldsymbol x-\boldsymbol\mu)]}{(2\pi)^{n/2}\sqrt{\det C}},\quad C>0.\]
\[Y=GX+b:\quad \mu_Y=G\mu_X+b,\quad C_Y=GC_XG^T.\]
联合高斯时，零互协方差等价于对应分量独立；只知道各边缘为高斯不够。
\key{零均值联合高斯的四阶矩（由特征函数或矩母函数导出）}
\[\E[X_1X_2X_3X_4]=C_{12}C_{34}+C_{13}C_{24}+C_{14}C_{23}.\]
\[Y=aX^2,\quad R_Y(s,t)=a^2\{C_X(s,s)C_X(t,t)+2C_X(s,t)^2\}.\]
\key{独立均匀相位与积分}
\[\E\cos(\alpha+\Theta)=0,\quad \E[\cos(\alpha+\Theta)\cos(\beta+\Theta)]=\tfrac12\cos(\alpha-\beta).\]
\[Y(t)=\int_0^tX(u)du:\quad R_Y(s,t)=\int_0^s\int_0^t R_X(u,v)\,dv\,du.\]
最后一式需要相应可积性以交换期望和积分；非平稳过程不能擅自把$R_X(u,v)$写成仅依赖时差。
\pagehead{13\quad 来源、验证与局限}
\key{材料位置（工作区相对路径）}
\begin{bul}
\item 教材：output/教材/SSA\_Textbook\_CH2\_zh.tex，书pp.37--62，习题2.1--2.21。
\item 作业：output/作业/SSA\_Homework\_2\_English.tex及英文详解；必做与选做清单已逐题对应。
\item 真题：试卷/下四份文件，文件尾缀分别为202212、202312、202412、202512\_1；学年和题号对应第3--6节审计表。
\item 历史基准：output/试卷/SSA\_CH1\_Exam\_Frequency\_Report.tex，保留其第一章主归属。
\item 方法：.dsh/skills/exam-frequency-report/SKILL.md及references/course-materials.md。
\item 本章机器可读审计：同目录SSA\_CH2\_Exam\_Tagging.json。
\end{bul}
\key{核验方式}
四套卷共16页全部渲染核对，题目页检查公式、上下标、分值和图中读数；文本抽取只作定位。
统计由同一份逐小问数据生成，检验每卷100分、本章36分、第一章既有总数不变。
配套题面、解答按同一题序生成；每题独立页；源码自包含；编译日志检查溢出和缺字，并检查PDF渲染。
\key{边界与局限}
\begin{bul}
\item 只有4套卷；本报告给出样本内描述，不能推断下一次试卷的概率或保证。
\item 归属是教学分类，不是官方阅卷细则；15分高斯四阶矩等边界已明示。如将来统一重分类，必须同步调整相邻章，不能单方面重复加分。
\item 间接标签只列解题关键知识，不把每道算期望的题都泛标为第二章；表内零涉及不等于知识毫无关系。
\item 原题条件不足或协方差奇异时，答案给出数学上成立的结论与条件式答案，不为了写出密度而默认增加条件。
\end{bul}
\vfill\textcolor{accentD}{编制日期：2026-10-13。三件套为自制学习材料。}
\end{document}
